\documentclass[9pt,a4paper,twocolumn]{ctexart}

% --- Packages ---
\usepackage[top=1.2cm, bottom=1.2cm, left=1.2cm, right=1.2cm, columnsep=0.8cm]{geometry}
\usepackage{hyperref}
\usepackage{graphicx}
\usepackage{float}
\usepackage{longtable}
\usepackage{tabularx}
\usepackage{booktabs}
\usepackage{enumitem}
\usepackage{xcolor}
\usepackage{fancyhdr}
\usepackage{listings}
\usepackage{fontspec}
\usepackage{titlesec}
\usepackage{caption}

% --- Code block styling ---
\definecolor{codebg}{RGB}{245,245,245}
\definecolor{codeframe}{RGB}{220,220,220}
\definecolor{codekw}{RGB}{0,0,192}
\definecolor{codecomment}{RGB}{0,128,0}
\definecolor{codestring}{RGB}{163,21,21}
\definecolor{codenum}{RGB}{128,0,128}
\definecolor{codepunct}{RGB}{90,90,90}
\definecolor{badgegreen}{RGB}{34,139,34}
\definecolor{badgeblue}{RGB}{30,144,255}
\definecolor{badgeblack}{RGB}{50,50,50}

% --- Difficulty badges (rounded corners via tikz, pre-rendered in saveboxes) ---
\usepackage{tikz}
\newsavebox{\badgechubox}
\savebox{\badgechubox}{\tikz[baseline]{\node[fill=badgegreen,text=white,rounded corners=2.5pt,inner xsep=4pt,inner ysep=1.5pt,font=\sffamily\footnotesize\bfseries]{初级};}}
\newsavebox{\badgezhongbox}
\savebox{\badgezhongbox}{\tikz[baseline]{\node[fill=badgeblue,text=white,rounded corners=2.5pt,inner xsep=4pt,inner ysep=1.5pt,font=\sffamily\footnotesize\bfseries]{中级};}}
\newsavebox{\badgegaobox}
\savebox{\badgegaobox}{\tikz[baseline]{\node[fill=badgeblack,text=white,rounded corners=2.5pt,inner xsep=4pt,inner ysep=1.5pt,font=\sffamily\footnotesize\bfseries]{高级};}}
\newcommand{\lvchu}{\usebox{\badgechubox}}
\newcommand{\lvzhong}{\usebox{\badgezhongbox}}
\newcommand{\lvgao}{\usebox{\badgegaobox}}
\pdfstringdefDisableCommands{%
  \def\lvchu{[初级]}%
  \def\lvzhong{[中级]}%
  \def\lvgao{[高级]}%
}

\lstset{
  basicstyle=\ttfamily\scriptsize,
  keywordstyle=\color{codekw}\bfseries,
  commentstyle=\color{codecomment}\itshape,
  stringstyle=\color{codestring},
  backgroundcolor=\color{codebg},
  frame=single,
  rulecolor=\color{codeframe},
  breaklines=true,
  breakatwhitespace=false,
  breakindent=0pt,
  postbreak=\mbox{\textcolor{red}{$\hookrightarrow$}\space},
  numbers=left,
  numberstyle=\tiny\color{gray},
  columns=flexible,
  keepspaces=true,
  showstringspaces=false,
  tabsize=2,
  aboveskip=4pt,
  belowskip=4pt,
  extendedchars=true,
  literate={，}{{，}}1 {；}{{；}}1 {！}{{！}}1 {？}{{？}}1 {“}{{“}}1 {”}{{”}}1 {（}{{（}}1 {）}{{）}}1 {《}{{《}}1 {》}{{》}}1,
}

% --- Extra language definitions (no built-in listings support) ---
\lstdefinelanguage{json}{
  morekeywords={true,false,null},
  sensitive=true,
  morestring=[b]",
  comment=[l]{//},
  morecomment=[s]{/*}{*/},
}
\lstdefinelanguage{yaml}{
  morekeywords={true,false,null,yes,no,on,off},
  sensitive=false,
  morecomment=[l]{\#},
  morestring=[b]",
  morestring=[d]',
}
\lstdefinelanguage{http}{
  morekeywords={GET,POST,PUT,DELETE,PATCH,HEAD,OPTIONS,HTTP,Host,Accept,Authorization,Content-Type,Content-Length,User-Agent},
  sensitive=true,
  morecomment=[l]{\#},
  morestring=[b]",
}
\lstdefinelanguage{TypeScript}{
  morekeywords={abstract,any,as,async,await,break,case,catch,class,const,continue,debugger,declare,default,delete,do,else,enum,export,extends,false,finally,for,from,function,get,if,implements,import,in,instanceof,interface,keyof,let,module,namespace,never,new,null,number,of,package,private,protected,public,readonly,return,set,static,string,super,switch,symbol,this,throw,true,try,type,typeof,undefined,unknown,var,void,while,with,yield,Record,Array,Promise,AsyncIterable,ReadableStream,Date,Error,Object,JSON},
  sensitive=true,
  morecomment=[l]{//},
  morecomment=[s]{/*}{*/},
  morestring=[b]",
  morestring=[b]',
  morestring=[b]`{`},
}

% --- Compact spacing ---
\linespread{0.95}
\setlength{\parskip}{1pt plus 1pt minus 1pt}
\setlength{\parindent}{0pt}

% --- Table styling ---
\renewcommand{\arraystretch}{1.1}

% --- Heading styling (numbered) ---
\titleformat{\section}{\normalsize\bfseries}{\thesection\quad}{0em}{}
\titlespacing{\section}{0pt}{6pt}{2pt}
\titleformat{\subsection}{\small\bfseries}{\thesubsection\quad}{0em}{}
\titlespacing{\subsection}{0pt}{4pt}{1pt}
\titleformat{\subsubsection}{\small\bfseries}{\thesubsubsection\quad}{0em}{}
\titlespacing{\subsubsection}{0pt}{3pt}{1pt}

% --- Page style ---
\pagestyle{fancy}
\fancyhf{}
\fancyhead[L]{\small\emph{\leftmark}}
\fancyhead[R]{\small\thepage}
\renewcommand{\headrulewidth}{0.4pt}
\renewcommand{\sectionmark}[1]{}  % prevent sections from overwriting leftmark

% --- Hyperref setup ---
\hypersetup{
  colorlinks=true,
  linkcolor=blue,
  urlcolor=blue,
  citecolor=blue,
}

% --- Caption format ---
\DeclareCaptionFormat{myformat}{\small\bfseries #1#2#3}
\captionsetup{format=myformat}

\begin{document}

\title{RAG 高级检索：混合检索、Rerank 与 GraphRAG}
\date{}
\maketitle
\markboth{Python 版 Agent 知识}{ Python 版 Agent 知识}

\begin{quote}
语言策略：code（Java / Python / TypeScript）
\end{quote}



\section*{TL;DR}

\begin{enumerate}[itemsep=1pt, topsep=2pt]
  \item 基础 RAG 的主要瓶颈不是“有没有检索”，而是\textbf{召回不准、上下文杂乱、问题表达不稳定}。
  \item 生产上优先按顺序做三件事：\textbf{混合检索 + RRF 融合 → Rerank 精排 → 查询改写}；GraphRAG 只用于多跳关系问题。
  \item 每条检索链路都要可评测、可回放：命中率、MRR、引用准确率必须进入回归门禁。
\end{enumerate}


\section*{1. 从基础 RAG 到高级检索}

{\footnotesize
\begin{tabular}{|p{\dimexpr 0.143\columnwidth-2\tabcolsep\relax}|p{\dimexpr 0.476\columnwidth-2\tabcolsep\relax}|p{\dimexpr 0.381\columnwidth-2\tabcolsep\relax}|}
\hline
\textbf{阶段} & \textbf{做法} & \textbf{解决什么} \\
\hline
基础 RAG & 单路向量召回 top-k & 语义相似问题 \\
\hline
混合检索 & 向量 + BM25 + 元数据过滤 & 专有名词、编号、时间范围 \\
\hline
融合 & RRF 或加权融合多路结果 & 多路排序尺度不统一 \\
\hline
精排 & Cross-Encoder Rerank & 前 50 条里选最相关的 5 条 \\
\hline
查询改写 & 改写、扩展、HyDE & 口语化、多义词、查询过短 \\
\hline
图检索 & GraphRAG & 多跳关系与全局总结 \\
\hline
\end{tabular}
}



\section*{2. Chunk 策略}

{\footnotesize
\begin{tabular}{|p{\dimexpr 0.258\columnwidth-2\tabcolsep\relax}|p{\dimexpr 0.226\columnwidth-2\tabcolsep\relax}|p{\dimexpr 0.516\columnwidth-2\tabcolsep\relax}|}
\hline
\textbf{策略} & \textbf{适用场景} & \textbf{注意} \\
\hline
固定长度 & 快速起步 & 语义会被切断 \\
\hline
段落/标题 & 结构化文档 & 需要保留层级元数据 \\
\hline
句子窗口 & 精读任务 & 检索用小句，送入模型带上下文窗口 \\
\hline
父子文档 & 命中后扩上下文 & 父块要单独存储 \\
\hline
表格/代码独立切 & 技术文档 & 不要按行切开 \\
\hline
\end{tabular}
}



原则：\textbf{切分结构要与检索结构一致}；每次切分策略变更都要重跑评测集。



\section*{3. 查询处理}

\begin{itemize}[itemsep=1pt, topsep=2pt]
  \item 查询改写：把“上次那个报错怎么修”改写成“订单服务 OOM 报错修复方法”；
  \item 查询扩展：生成同义词或多个子查询；
  \item HyDE：先用 LLM 生成假设答案，再用答案向量去召回；
  \item 多轮对话：先做指代消解，把“它”替换为具体对象，再检索；
  \item 元数据过滤：时间、租户、文档类型先过滤，再做语义召回。
\end{itemize}


\section*{4. 混合检索与 RRF}


BM25 擅长专有名词和编号，向量检索擅长语义。两路分数不可直接相加，常用 Reciprocal Rank Fusion：



\begin{lstlisting}
RRF_score(doc) = sum( 1 / (60 + rank_i(doc)) )
\end{lstlisting}



实现要点：先分别取每路 top-100，再融合取 top-20，最后交给 Rerank 取 top-5。



\section*{5. Rerank}

\begin{itemize}[itemsep=1pt, topsep=2pt]
  \item Cross-Encoder 比 Bi-Encoder 精度高、延迟高，所以只对候选集精排；
  \item 离线评测用 MRR、nDCG 和引用准确率，不要只看“感觉更相关”；
  \item 候选集通常 20–50 条，精排后送入上下文 3–8 条；
  \item 工具选择、模型切换必须过回归门禁。
\end{itemize}


\section*{6. GraphRAG 什么时候上}

{\footnotesize
\begin{tabular}{|p{\dimexpr 0.333\columnwidth-2\tabcolsep\relax}|p{\dimexpr 0.286\columnwidth-2\tabcolsep\relax}|p{\dimexpr 0.381\columnwidth-2\tabcolsep\relax}|}
\hline
\textbf{问题类型} & \textbf{普通 RAG} & \textbf{GraphRAG} \\
\hline
单文档问答 & 够用 & 不需要 \\
\hline
专有名词/编号 & 混合检索 & 不需要 \\
\hline
多跳关系 & 容易断链 & 更适合 \\
\hline
全局总结 & 容易偏 & 社区摘要更有结构 \\
\hline
\end{tabular}
}



GraphRAG 成本高、更新难，先确认任务真的需要实体关系和社区摘要，再引入。



\section*{7. 检索评测}

\begin{itemize}[itemsep=1pt, topsep=2pt]
  \item 指标：命中率、MRR、nDCG、上下文引用准确率；
  \item 评测集：真实问题 + 期望文档 ID + 禁止返回文档；
  \item 每次改动 Chunk、Embedding、融合参数、Rerank 模型都必须回归；
  \item 线上 Bad Case 回流：检索失败、引用错误要进入下一轮评测集。
\end{itemize}


\section*{8. Python 版检索融合实现}


\begin{lstlisting}[language=Python]
from collections import defaultdict

def rrf_fusion(ranked_lists: list[list[tuple[str, float, str]]],
               k: float = 60.0, top_n: int = 20) -> list[tuple[str, float, str]]:
    """ranked_lists: 每路召回结果 [(doc_id, score, source), ...]，按分数降序。"""
    scores: dict[str, float] = defaultdict(float)
    sources: dict[str, str] = {}
    for ranked in ranked_lists:
        for i, (doc_id, _score, source) in enumerate(ranked):
            scores[doc_id] += 1.0 / (k + i + 1)
            sources.setdefault(doc_id, source)
    fused = sorted(scores.items(), key=lambda item: item[1], reverse=True)[:top_n]
    return [(doc_id, score, sources[doc_id]) for doc_id, score in fused]
\end{lstlisting}



\section*{9. 延伸阅读}

\begin{itemize}[itemsep=1pt, topsep=2pt]
  \item \href{../10-面试/02-RAG面试题.md}{RAG 面试题}
  \item \href{../05-评测与质量/01-Agent评测体系.md}{Agent 评测体系}
  \item \href{../07-系统设计/01-AI应用系统设计.md}{AI 应用系统设计}
\end{itemize}

\end{document}
